% !TEX program = lualatex
% !TEX options = --shell-escape
\documentclass[a4paper,12pt]{article}

% ============================================================
% preamble.tex — Package Loading & Global Configuration
% Compile with: lualatex --shell-escape main.tex
% ============================================================

% ====== FONT CONFIG ======
\newcommand{\MainFont}{Optima}     % LaTeX default (Computer Modern successor)
\newcommand{\SansFont}{Latin Modern Sans}      % LaTeX default sans
\newcommand{\MonoFont}{JetBrainsMono Nerd Font Mono}
\newcommand{\CJKMainFont}{jf-jinxuan-Book}
\newcommand{\CJKSansFont}{jf-jinxuanlatte-2.0}
\newcommand{\CJKMonoFont}{JetBrainsMono Nerd Font Mono}

% ====== PAGE LAYOUT ======
\newcommand{\PaperSize}{a4paper}
\newcommand{\MarginLeft}{1.5cm}
\newcommand{\MarginRight}{1.5cm}
\newcommand{\MarginTop}{2cm}
\newcommand{\MarginBottom}{2cm}
\newcommand{\BaseLineStretch}{1.15}   % line spacing multiplier

% ====== COLORS ======
\newcommand{\TitleColor}{black}
\newcommand{\HeadingColor}{black}
\newcommand{\BodyColor}{black}
\newcommand{\LinkColor}{MidnightBlue}
\newcommand{\CaptionColor}{black}
\newcommand{\CodeBgColor}{gray!10}

% ============================================================
% PACKAGE LOADING
% ============================================================

% --- Typography ---
\usepackage{microtype}                     % typographic micro-adjustments

% --- Fonts ---
\usepackage{fontspec}                      % system font loading (LuaLaTeX)
\usepackage{luatexja-fontspec}             % CJK support for LuaLaTeX

% --- Page Layout ---
\usepackage{geometry}                      % page margins
\usepackage{fancyhdr}                      % header / footer
\usepackage{multicol}                      % multi-column layouts
\usepackage{marginnote}                    % margin notes

% --- Colors & Graphics ---
\usepackage[dvipsnames, svgnames]{xcolor}  % extended color names
\usepackage{graphicx}                      % image inclusion
\graphicspath{{figures/}{assets/}}

% --- Paragraph Spacing ---
% parskip package: sets \parskip and \parindent, and fixes interactions
% with display math, lists, and toc (prevents extra spacing).
% Replaces manual \setlength{\parindent}{...} and \setlength{\parskip}{...}.
% To adjust: \usepackage[skip=<size> plus <stretch>, indent=<size>]{parskip}
\usepackage[skip=6pt plus 2pt]{parskip}

% --- Math ---
\usepackage{amsmath, amssymb, amsthm}

% --- Tables ---
\usepackage{booktabs}                      % professional tables

% --- Lists ---
\usepackage{enumitem}                      % customizable lists

% --- Code Highlighting ---
\usepackage{listings}                      % syntax highlighting (no external deps)

% --- Section Formatting ---
\usepackage{titlesec}

% --- Captions ---
\usepackage{caption}

% --- Boxes ---
\usepackage[most]{tcolorbox}               % colored boxes / callouts
\usepackage{tikz}                          % inline graphics (link badges, etc.)
\usetikzlibrary{shadows,backgrounds,fit,positioning}

% --- Algorithms ---
\usepackage[tworuled, vlined, linesnumbered]{algorithm2e}

% --- Bibliography ---
\usepackage[
    backend=biber,
    style=numeric,
    sorting=none
]{biblatex}

% --- Todo Notes ---
\setlength{\marginparwidth}{2cm}
\usepackage{todonotes}

% --- Strikethrough ---
\usepackage[normalem]{ulem}                % \sout{text} for strikethrough, normalem keeps \emph italic

% --- Hyperlinks (must be near last) ---
\usepackage{hyperref}

% --- Smart Cross-References (must be after hyperref) ---
\usepackage{cleveref}                      % \cref{fig:xxx} → "Figure 1" automatically

% ============================================================
% PACKAGE CONFIGURATION (uses variables defined above)
% ============================================================

% --- Fonts ---
\setmainfont[Ligatures=TeX]{\MainFont}
\setsansfont{\SansFont}
\setmonofont[Scale=0.85]{\MonoFont}

\setmainjfont[BoldFont=jf-jinxuan-Bold]{\CJKMainFont}
\setsansjfont{\CJKSansFont}

% --- Page Geometry ---
\geometry{
    \PaperSize,
    left=\MarginLeft,
    right=\MarginRight,
    top=\MarginTop,
    bottom=\MarginBottom,
    headheight=14pt,
    headsep=8pt,
    footskip=20pt,
}

% --- Paragraph ---
\renewcommand{\baselinestretch}{\BaseLineStretch}

% --- Display Math Spacing ---
% \normalsize resets \abovedisplayskip to ~14pt on every call;
% append our tighter values so they persist after font-size changes.
\appto\normalsize{%
  \abovedisplayskip=4pt plus 2pt minus 2pt
  \belowdisplayskip=4pt plus 2pt minus 2pt
  \abovedisplayshortskip=0pt plus 2pt
  \belowdisplayshortskip=2pt plus 2pt minus 2pt
}
\normalsize

% --- Body Color ---
\color{\BodyColor}

% --- Header / Footer (base setup; styles configure content) ---
\pagestyle{fancy}
\fancyhf{}

% --- Hyperlinks ---
\hypersetup{
    colorlinks=true,
    linkcolor=\LinkColor,
    filecolor=magenta,
    urlcolor=\LinkColor,
    citecolor=\LinkColor,
}

% --- Listings (code highlighting — light theme, rendered inside tcolorbox) ---
\definecolor{CodeBg}{HTML}{f6f8fa}
\definecolor{CodeFg}{HTML}{24292e}
\definecolor{CodeKeyword}{HTML}{d73a49}
\definecolor{CodeString}{HTML}{032f62}
\definecolor{CodeComment}{HTML}{6a737d}
\definecolor{CodeNumber}{HTML}{a0a8b2}
\definecolor{CodeFunction}{HTML}{6f42c1}
\definecolor{CodeAccent}{HTML}{d0d7de}
\lstset{
    basicstyle=\small\ttfamily\color{CodeFg},
    numbers=left,
    numberstyle=\tiny\ttfamily\color{CodeNumber},
    breaklines=true,
    tabsize=4,
    frame=none,
    showstringspaces=false,
    keywordstyle=\color{CodeKeyword}\bfseries,
    commentstyle=\color{CodeComment}\itshape,
    stringstyle=\color{CodeString},
    columns=fullflexible,
    keepspaces=true,
    xleftmargin=14pt,
}

% --- Captions ---
\captionsetup{
    font=small,
    labelfont={bf},
    textfont={color=\CaptionColor},
    skip=6pt,
}

% --- Algorithm2e ---
\SetAlgoLined
\SetKwInput{KwInput}{Input}
\SetKwInput{KwOutput}{Output}
\SetNlSty{textnormal}{\color{gray!40}}{}  % light gray line numbers
% Separator between Input/Output and algorithm body (matches ruled line style)
\newcommand{\AlgoSep}{%
  \vspace{2pt}%
  \nointerlineskip
  \hbox to 0pt{\vrule width\linewidth height\algoheightrule depth0pt\hss}%
  \nointerlineskip
  \vspace{\dimexpr4pt+\algoheightrule}%
}

% --- Lists ---
\setlist{nosep, leftmargin=1.5em}

% ============================================================
% macros.tex — Style System & Custom Commands
% ============================================================

% ======================================================================
% A. STYLE SYSTEM
% ======================================================================

% --- Document info setters ---
\let\DocTitle\empty
\let\DocSubtitle\empty
\let\DocContact\empty
\let\DocAuthor\empty
\let\DocDate\empty
\let\DocHeader\empty

\newcommand{\doctitle}[1]{\def\DocTitle{#1}}
\newcommand{\docsubtitle}[1]{\def\DocSubtitle{#1}}
\newcommand{\doccontact}[1]{\def\DocContact{#1}}
\newcommand{\docauthor}[1]{\def\DocAuthor{#1}}
\newcommand{\docdate}[1]{\def\DocDate{#1}}
\newcommand{\docheader}[1]{\def\DocHeader{#1}}

\makeatletter

% --- Style dispatcher ---
\newcommand{\usestyle}[1]{%
    \ifcsname style@#1\endcsname
        \csname style@#1\endcsname
    \else
        \PackageError{macros}{Unknown style '#1'}{%
            Available styles: academic, homework, note, minimal, custom}%
    \fi
}

% ------------------------------------------------------------------
% Style: academic
%   Title: centered, large title + subtitle + author + date
%   Sections: \Large bold with underline, numbered
%   Header: empty | Footer: centered page number
% ------------------------------------------------------------------
\newcommand{\style@academic}{%
    % Title
    \renewcommand{\maketitle}{%
        \begin{center}
            {\fontsize{30pt}{36pt}\selectfont\bfseries\color{\TitleColor}\DocTitle\par}
            \vspace{6pt}
            \ifx\DocSubtitle\empty\else
                {\fontsize{12pt}{16pt}\selectfont\color{\HeadingColor}\DocSubtitle\par}
                \vspace{4pt}
            \fi
            \ifx\DocAuthor\empty\else
                {\normalsize\DocAuthor\par}
                \vspace{2pt}
            \fi
            \ifx\DocDate\empty\else
                {\small\color{gray}\DocDate\par}
                \vspace{2pt}
            \fi
            \ifx\DocContact\empty\else
                {\small\color{gray}\DocContact\par}
            \fi
        \end{center}
        \vspace{12pt}
    }%
    % Sections
    \titleformat{\section}
        {\color{\HeadingColor}\Large\bfseries}
        {\thesection}{1em}{}
        [\vspace{-4pt}\color{\HeadingColor}\titlerule]
    \titlespacing{\section}{0pt}{12pt}{2pt}
    \titleformat{\subsection}
        {\color{\HeadingColor}\large\bfseries}
        {\thesubsection}{1em}{}
    \titlespacing{\subsection}{0pt}{10pt}{0pt}
    \titleformat{\subsubsection}
        {\color{\HeadingColor}\normalsize\bfseries}
        {\thesubsubsection}{1em}{}
    \titlespacing{\subsubsection}{0pt}{6pt}{0pt}
    % Header / Footer
    \fancyhf{}
    \fancyfoot[C]{\thepage}
    \renewcommand{\headrulewidth}{0pt}
    \renewcommand{\footrulewidth}{0pt}
}

% ------------------------------------------------------------------
% Style: homework
%   Title: left-right aligned table (title+subtitle left, date+author right), hrule
%   Sections: \large bold, no underline, NO numbering (type anything in heading)
%   Header: left=\docheader (defaults to subtitle), right=author | Footer: centered page number
%   Header rule: 0.4pt
% ------------------------------------------------------------------
\newcommand{\style@homework}{%
    % Title (first page has no header)
    \renewcommand{\maketitle}{%
        \thispagestyle{plain}%
        \noindent
        \begin{tabular*}{\textwidth}{@{}l@{\extracolsep{\fill}}r@{}}
            {\fontsize{30pt}{36pt}\selectfont\bfseries\color{\TitleColor}\DocTitle}
            & {\normalsize\DocDate} \\[4pt]
            {\fontsize{12pt}{16pt}\selectfont\color{\HeadingColor}\DocSubtitle}
            & {\normalsize\DocAuthor} \\
        \end{tabular*}
        \ifx\DocContact\empty\else
            \\[2pt]\noindent{\small\color{gray}\DocContact}
        \fi
        \vspace{4pt}
        \hrule
        \vspace{12pt}
    }%
    % Sections: all unnumbered (type anything in heading); h1 has rule
    \titleformat{\section}
        {\color{\HeadingColor}\LARGE\bfseries}
        {}{0em}{}
        [\vspace{-4pt}\color{gray!70}\titlerule]
    \titlespacing{\section}{0pt}{10pt}{2pt}
    \titleformat{\subsection}
        {\color{\HeadingColor}\Large\bfseries}
        {}{0em}{}
    \titlespacing{\subsection}{0pt}{8pt}{-2pt}
    \titleformat{\subsubsection}
        {\color{\HeadingColor}\large\bfseries}
        {}{0em}{}
    \titlespacing{\subsubsection}{0pt}{6pt}{-2pt}
    % Header / Footer
    \fancyhf{}
    \fancyhead[L]{\small\ifx\DocHeader\empty\DocSubtitle\else\DocHeader\fi}
    \fancyhead[R]{\small\DocAuthor}
    \fancyfoot[C]{\small\thepage}
    \renewcommand{\headrulewidth}{0.4pt}
    \renewcommand{\footrulewidth}{0pt}
    \fancypagestyle{plain}{\fancyhf{}\fancyfoot[C]{\small\thepage}\renewcommand{\headrulewidth}{0pt}}%
}

% ------------------------------------------------------------------
% Style: note
%   Title: left-aligned 24pt title, subtitle, author · date inline, hrule
%   Sections: \LARGE bold with underline, numbered
%   Header: left=\docheader (defaults to title), right=author | Footer: centered page number
%   Header rule: 0.4pt
% ------------------------------------------------------------------
\newcommand{\style@note}{%
    % Title (first page has no header)
    \renewcommand{\maketitle}{%
        \thispagestyle{plain}%
        \noindent
        {\fontsize{24pt}{30pt}\selectfont\bfseries\color{\TitleColor}\DocTitle\par}
        \vspace{4pt}
        \ifx\DocSubtitle\empty\else
            {\vspace{-6pt}\fontsize{12pt}{16pt}\selectfont\color{\HeadingColor}\DocSubtitle\par}
        \fi
        \ifx\DocAuthor\empty\else
            {\vspace{-6pt}\normalsize\DocAuthor}%
            \ifx\DocDate\empty\else{ · \DocDate}\fi
            \par
        \fi
        \ifx\DocContact\empty\else
            {\small\color{gray}\DocContact\par}
        \fi
        \vspace{4pt}
        \hrule
        \vspace{12pt}
    }%
    % Sections: all unnumbered (type anything in heading); h1 has rule
    \titleformat{\section}
        {\color{\HeadingColor}\LARGE\bfseries}
        {\thesection.}{0.25em}{}
        [\vspace{-4pt}\color{gray!70}\titlerule]
    \titlespacing{\section}{0pt}{10pt}{2pt}
    \titleformat{\subsection}
        {\color{\HeadingColor}\Large\bfseries}
        {}{0em}{}
    \titlespacing{\subsection}{0pt}{8pt}{-2pt}
    \titleformat{\subsubsection}
        {\color{\HeadingColor}\large\bfseries}
        {}{0em}{}
    \titlespacing{\subsubsection}{0pt}{6pt}{-2pt}
    % Header / Footer
    \fancyhf{}
    \fancyhead[L]{\small\ifx\DocHeader\empty\DocTitle\else\DocHeader\fi}
    \fancyhead[R]{\small\DocAuthor}
    \fancyfoot[C]{\small\thepage}
    \renewcommand{\headrulewidth}{0.4pt}
    \renewcommand{\footrulewidth}{0pt}
    \fancypagestyle{plain}{\fancyhf{}\fancyfoot[C]{\small\thepage}\renewcommand{\headrulewidth}{0pt}}%
}

% ------------------------------------------------------------------
% Style: minimal
%   Title: left-aligned, subtitle · author · date inline
%   Sections: \large bold, no underline, NO numbering
%   Header: empty | Footer: centered page number
% ------------------------------------------------------------------
\newcommand{\style@minimal}{%
    % Title
    \renewcommand{\maketitle}{%
        \noindent
        {\fontsize{30pt}{36pt}\selectfont\bfseries\color{\TitleColor}\DocTitle\par}
        \vspace{4pt}
        \ifx\DocSubtitle\empty\else
            {\fontsize{12pt}{16pt}\selectfont\color{\HeadingColor}\DocSubtitle\par}
            \vspace{2pt}
        \fi
        \ifx\DocAuthor\empty\else
            {\normalsize\DocAuthor}%
            \ifx\DocDate\empty\else{ · \DocDate}\fi
            \par
        \fi
        \vspace{12pt}
    }%
    % Sections (no numbering)
    \titleformat{\section}
        {\color{\HeadingColor}\Large\bfseries}
        {}{0em}{}
    \titlespacing{\section}{0pt}{12pt}{2pt}
    \titleformat{\subsection}
        {\color{\HeadingColor}\large\bfseries}
        {}{0em}{}
    \titlespacing{\subsection}{0pt}{8pt}{0pt}
    \titleformat{\subsubsection}
        {\color{\HeadingColor}\normalsize\bfseries}
        {}{0em}{}
    \titlespacing{\subsubsection}{0pt}{6pt}{0pt}
    % Header / Footer
    \fancyhf{}
    \fancyfoot[C]{\thepage}
    \renewcommand{\headrulewidth}{0pt}
    \renewcommand{\footrulewidth}{0pt}
}

% ------------------------------------------------------------------
% Style: custom (academic as base; override with \renewcommand / \titleformat after)
% ------------------------------------------------------------------
\newcommand{\style@custom}{%
    \style@academic
}

\makeatother

% ======================================================================
% B. CODE BLOCKS
% ======================================================================

% codeblock environment: \begin{codeblock}[python] ... \end{codeblock}
% Light-themed code block with left accent border (uses tcolorbox + listings)
\newtcblisting{codeblock}[1][python]{
    colback=CodeBg,
    colframe=CodeBg,
    listing only,
    listing options={language=#1},
    arc=2pt,
    boxrule=0pt,
    borderline west={2.5pt}{0pt}{CodeAccent},
    left=8pt, right=8pt, top=4pt, bottom=4pt,
    before skip=8pt,
    after skip=8pt,
}

% inline code — rounded background pill
\newtcbox{\inlinecode}{
    on line,
    colback=\CodeBgColor,
    colframe=\CodeBgColor,
    boxrule=0pt,
    arc=3pt,
    boxsep=0pt,
    left=4pt, right=4pt, top=1.5pt, bottom=1.5pt,
    fontupper=\small\ttfamily,
}

% ======================================================================
% D. CALLOUT / NOTE BOXES
% ======================================================================

% Usage: \begin{bluebox}...\end{bluebox}
%        \begin{bluebox}[Custom Title]...\end{bluebox}

\definecolor{BoxBlue}{HTML}{4a90d9}
\definecolor{BoxAmber}{HTML}{e6a117}
\definecolor{BoxGreen}{HTML}{3a9a5c}
\definecolor{BoxRed}{HTML}{d94452}

\newtcolorbox{bluebox}[1][Note]{
    colback=BoxBlue!6, colframe=BoxBlue, coltitle=white,
    title={\smash{\textbf{#1}}\vphantom{N}}, fonttitle=\sffamily,
    boxrule=0.5pt, arc=2pt,
    left=6pt, right=6pt, top=4pt, bottom=4pt,
}

\newtcolorbox{amberbox}[1][Warning]{
    colback=BoxAmber!8, colframe=BoxAmber, coltitle=white,
    title={\smash{\textbf{#1}}\vphantom{N}}, fonttitle=\sffamily,
    boxrule=0.5pt, arc=2pt,
    left=6pt, right=6pt, top=4pt, bottom=4pt,
}

\newtcolorbox{greenbox}[1][Tip]{
    colback=BoxGreen!6, colframe=BoxGreen, coltitle=white,
    title={\smash{\textbf{#1}}\vphantom{N}}, fonttitle=\sffamily,
    boxrule=0.5pt, arc=2pt,
    left=6pt, right=6pt, top=4pt, bottom=4pt,
}

\newtcolorbox{redbox}[1][Danger]{
    colback=BoxRed!6, colframe=BoxRed, coltitle=white,
    title={\smash{\textbf{#1}}\vphantom{N}}, fonttitle=\sffamily,
    boxrule=0.5pt, arc=2pt,
    left=6pt, right=6pt, top=4pt, bottom=4pt,
}

% ======================================================================
% E. THEOREM ENVIRONMENTS
% ======================================================================

\newtheorem{theorem}{Theorem}[section]
\newtheorem{definition}{Definition}[section]
\newtheorem{lemma}{Lemma}[section]
\newtheorem{corollary}{Corollary}[section]
\newtheorem{proposition}{Proposition}[section]

% ======================================================================
% F. ALGORITHM (algorithm2e defaults set in preamble)
% ======================================================================

% ======================================================================
% G. TODO
% ======================================================================

% Margin todo note
\newcommand{\todox}[1]{\todo[inline, color=yellow!40]{#1}}

% Checkbox-style todo
\newcommand{\todomark}[1]{%
    $\square$~{#1}%
}
\newcommand{\tododone}[1]{%
    $\boxtimes$~{\sout{#1}}%
}

% ======================================================================
% H. MULTI-COLUMN LAYOUTS
% ======================================================================

\newenvironment{twocol}{%
    \begin{multicols}{2}%
}{%
    \end{multicols}%
}

\newenvironment{threecol}{%
    \begin{multicols}{3}%
}{%
    \end{multicols}%
}

% ======================================================================
% I. BLOCKQUOTE
% ======================================================================

\newtcolorbox{blockquote}{
    blanker,
    borderline west={3pt}{0pt}{gray!60},
    left=12pt,
    top=4pt,
    bottom=4pt,
    before skip=8pt,
    after skip=8pt,
    fontupper=\itshape\color{gray!80!black},
}

% ======================================================================
% J. SHADOW BOX (generic wrapper: border + drop shadow)
% ======================================================================

% Usage: \shadowbox{\includegraphics[...]{...}}
%        \shadowbox{any content}
% Shadow is rendered via tikz [overlay] so it does NOT affect bounding box,
% meaning centering, spacing, and caption distance stay identical to bare content.
%
% --- CONFIG (tweak these) ---
\newcommand{\ShadowColor}{gray}        % shadow color
\newcommand{\ShadowBorderColor}{gray}  % border color
\newcommand{\ShadowBorderWidth}{0.4pt} % border line width
\newcommand{\ShadowLayers}{5}          % number of blur layers (more = smoother)
\newcommand{\ShadowXShift}{0.15mm}     % base X offset per layer (positive = right)
\newcommand{\ShadowYShift}{-0.15mm}    % base Y offset per layer (negative = down)
\newcommand{\ShadowSpread}{0.12mm}     % extra spread per layer (controls blur radius)
\newcommand{\ShadowBaseOpacity}{0.12}  % opacity of innermost shadow layer
\newcommand{\ShadowFadeRate}{0.018}    % opacity decrease per layer
% ---
\newcommand{\shadowbox}[1]{%
    \begin{tikzpicture}
        \node[inner sep=0pt] (C) {#1};
        \draw[\ShadowBorderColor, line width=\ShadowBorderWidth]
            (C.south west) rectangle (C.north east);
        \begin{scope}[on background layer, overlay]
            \foreach \s in {1,...,\ShadowLayers} {
                \pgfmathsetmacro{\op}{\ShadowBaseOpacity - \ShadowFadeRate*\s}
                \fill[\ShadowColor, opacity=\op]
                    ([shift={({\ShadowXShift+\ShadowSpread*\s},{\ShadowYShift-\ShadowSpread*\s})}]C.south west)
                    rectangle
                    ([shift={({\ShadowXShift+\ShadowSpread*\s},{\ShadowYShift-\ShadowSpread*\s})}]C.north east);
            }
        \end{scope}
    \end{tikzpicture}%
}

% ======================================================================
% K. MARGIN NOTE
% ======================================================================

\newcommand{\marginmark}[1]{%
    \marginnote{\footnotesize\color{gray}#1}%
}


% ====== DOCUMENT INFO ======
% 填入文件資訊，留空或不寫的欄位會自動跳過（不會出現空白行）
\doctitle{Template Reference}
\docsubtitle{A comprehensive guide to using the LaTeX template}   % 選填
\docauthor{Roger Fan}
\docdate{\today}
% \doccontact{email@example.com}   % 選填
% \docheader{Custom Header}       % 選填，homework 樣式的頁首左邊文字（預設為 subtitle）

% ====== STYLE ======
% 可用樣式：academic / homework / minimal / custom
%   academic — 標題置中，section 有底線，頁尾頁碼
%   homework — 標題左右對齊（標題+課名 左，日期+姓名 右），section 無編號，頁首左=\docheader（預設 subtitle）右=姓名
%   minimal  — 標題靠左，section 無編號，乾淨簡潔
%   custom   — 以 academic 為底，可在下方自行覆蓋
\usestyle{homework}

\begin{document}

\maketitle

% ============================================================
\section{Headings}
% ============================================================

% \section{標題}        — 第一層（homework 無編號，academic 自動編號 1. 2. 3.）
% \subsection{標題}     — 第二層
% \subsubsection{標題}  — 第三層
% 編號格式和是否顯示編號取決於 \usestyle 的設定

這是正文範例。Body text with both English and 中文混排。

\subsection{Heading 2 Example}

段落之間空一行就會自動產生間距（由 preamble.tex 的 ParSkip 控制）。

\subsubsection{Heading 3 Example}

第三層標題。

% ============================================================
\section{Text Formatting}
% ============================================================

% --- 基本格式 ---
\textbf{粗體 Bold}、
\textit{斜體 Italic}、
\texttt{等寬字 Monospace}、
\underline{底線 Underline}、
\sout{刪除線 Strikethrough}。

% --- 組合使用 ---
% 用大括號包住，裡面疊加效果
\textbf{\textit{粗斜體}}、
{\bfseries\color{red}粗體紅字}。

% --- 顏色 ---
% \textcolor{顏色}{文字}   — 指定文字變色
% 預設可用顏色：red, blue, green, orange, purple, gray, cyan, magenta, ...
% dvipsnames 額外顏色：MidnightBlue, ForestGreen, BrickRed, ...（見 xcolor 文件）
% 自訂顏色：\definecolor{MyColor}{HTML}{FF6B35}（在 preamble 或文件開頭定義）
\textcolor{red}{紅色文字}、
\textcolor{MidnightBlue}{深藍色文字}、
\textcolor{ForestGreen}{森林綠文字}。

% --- 字體大小（由小到大）---
% \tiny \scriptsize \footnotesize \small \normalsize \large \Large \LARGE \huge \Huge
% 用大括號限制範圍：{\large 這段比較大}，括號外恢復正常
{\footnotesize 小字}、
{\normalsize 正常}、
{\large 大字}、
{\Large 更大}。

% --- 特殊字元與空白 ---
% ~          — 不斷行空格（non-breaking space）
%              兩邊的字不會被拆到不同行，常用於 "Figure~1"、"Dr.~Smith"
% -          — 短橫線（hyphen），用於複合詞：well-known
% --         — 半形破折號（en dash），用於數字範圍：pp.~1--10、2020--2024
% ---        — 全形破折號（em dash），用於插入語：LaTeX --- a typesetting system --- is powerful.
% ``文字''   — 英文左右雙引號（鍵盤上 ` 在 Tab 上方，' 在 Enter 左邊）
% `文字'     — 英文左右單引號
% \%         — 顯示 % 符號（% 本身是註解符號）
% \&         — 顯示 & 符號（& 本身是表格欄位分隔符）
% \$         — 顯示 $ 符號（$ 本身是數學模式符號）
% \\         — 強制換行（不建議在正文中用，正文用空行分段落即可）

Hyphens and dashes: well-known, pp.~1--10, 2020--2024, LaTeX --- a powerful tool.

English quotes: ``double quotes'' and `single quotes'.

% --- 超連結 ---
\href{https://example.com}{可點擊的連結文字}。

% ============================================================
\section{Lists}
% ============================================================

% --- 無序列表 ---
% 預設符號為 •，可巢狀（符號自動變化）
\subsection{Unordered List}

\begin{itemize}
    \item First item
    \item Second item with sub-items:
    \begin{itemize}
        \item Sub-item A
        \item Sub-item B
    \end{itemize}
    \item Third item
\end{itemize}

% --- 有序列表 ---
\subsection{Ordered List}

\begin{enumerate}
    \item Step one
    \item Step two
    \item Step three
\end{enumerate}

% --- 自訂符號與起始編號 ---
% [label=\textbullet]   — 指定符號（\textbullet •, \textendash –, $\star$ ★）
% [start=5]             — 從 5 開始編號
% [label=(\alph*)]      — 改成 (a) (b) (c)
% [label=(\roman*)]     — 改成 (i) (ii) (iii)
\subsection{Custom Lists}

\begin{enumerate}[label=(\alph*)]
    \item First with letter label
    \item Second with letter label
\end{enumerate}

\begin{enumerate}[start=3]
    \item This starts at 3
    \item This is 4
\end{enumerate}

% --- 描述列表 ---
\subsection{Description List}

\begin{description}
    \item[Term A] Definition of term A.
    \item[Term B] Definition of term B.
\end{description}

% ============================================================
\section{Callout Boxes}
% ============================================================

% 四種顏色的提示框，每種都有預設標題，也可以自訂
% \begin{bluebox}              — 預設標題 "Note"
% \begin{bluebox}[Custom Title] — 自訂標題
% 同理：amberbox (Warning)、greenbox (Tip)、redbox (Danger)

\begin{bluebox}
This is an informational note with default title.
\end{bluebox}

\begin{amberbox}[Caution: Memory Leak]
This is a warning with a custom title.
\end{amberbox}

\begin{greenbox}
This is a tip with default title.
\end{greenbox}

\begin{redbox}[Critical Bug]
This is a danger alert with a custom title.
\end{redbox}

% ============================================================
\section{Blockquote}
% ============================================================

% 左側灰色豎線 + 斜體灰字，用於引言
\begin{blockquote}
The only way to do great work is to love what you do.

--- Steve Jobs
\end{blockquote}

% ============================================================
\section{Mathematics}
% ============================================================

% --- 行內數學 ---
% 用 $ ... $ 包住
Inline math: $E = mc^2$, $\sum_{i=1}^{n} i = \frac{n(n+1)}{2}$.

% --- 獨立公式（有編號）---
% equation 環境，自動編號，可加 \label 引用
\begin{equation}
    \int_{-\infty}^{\infty} e^{-x^2} \, dx = \sqrt{\pi}
    \label{eq:gaussian}
\end{equation}

Reference: \cref{eq:gaussian}.

% --- 獨立公式（無編號）---
% 用 \[ ... \] 或 equation* 環境
\[
    a^2 + b^2 = c^2
\]

% --- 多行對齊公式 ---
% align 環境，用 & 標記對齊點，用 \\ 換行
% 每行都有編號；不要編號的行在 \\ 前加 \nonumber
\begin{align}
    f(x) &= x^2 + 2x + 1     \label{eq:fx} \\
         &= (x + 1)^2         \nonumber
\end{align}

% --- 常用符號速查 ---
% \frac{分子}{分母}    — 分數
% \sqrt{x}, \sqrt[3]{x} — 平方根、立方根
% x^{2}, x_{i}         — 上標、下標
% \sum, \prod, \int     — 求和、連乘、積分
% \alpha \beta \gamma   — 希臘字母
% \leq \geq \neq       — ≤ ≥ ≠
% \in \notin \subset    — ∈ ∉ ⊂
% \to \Rightarrow       — → ⇒
% \infty                — ∞
% \cdot \times          — · ×
% \mathbb{R}            — ℝ（需 amssymb）
% \text{文字}           — 在數學模式中插入普通文字

% --- 定理環境 ---
% 可用：theorem, definition, lemma, corollary, proposition
% 自動編號，格式為 "Theorem 7.1"（section.序號）

\begin{theorem}
    For any triangle with sides $a$, $b$, and $c$, and the angle $C$ opposite side $c$:
    \[
        c^2 = a^2 + b^2 - 2ab\cos C
    \]
\end{theorem}

\begin{definition}
    A \textbf{group} $(G, \cdot)$ is a set $G$ with a binary operation $\cdot$ satisfying closure, associativity, identity, and invertibility.
\end{definition}

\begin{lemma}
    If $f$ is continuous on $[a, b]$, then $f$ is bounded on $[a, b]$.
\end{lemma}

% ============================================================
\section{Code}
% ============================================================

% --- 行內程式碼 ---
% \inlinecode{程式碼} — 灰底圓角小藥丸樣式
Inline code: \inlinecode{print("hello")} is a Python statement.

% --- 程式碼區塊 ---
% \begin{codeblock}[語言] ... \end{codeblock}
% 預設語言是 python，可改成其他 listings 支援的語言：
%   python, java, c, cpp, javascript, bash, sql, html, ruby, go, rust, ...
% 完整列表見 listings 套件文件

\begin{codeblock}[python]
def fibonacci(n: int) -> int:
    """Return the n-th Fibonacci number."""
    if n <= 1:
        return n
    return fibonacci(n - 1) + fibonacci(n - 2)

print(fibonacci(10))  # 55
\end{codeblock}

\begin{codeblock}[bash]
# shell script example
for i in $(seq 1 5); do
    echo "iteration $i"
done
\end{codeblock}

% ============================================================
\section{Figures}
% ============================================================

% figure 語法解析：
%   \begin{figure}[htbp]          — 浮動環境（LaTeX 自動決定最佳位置）
%     位置偏好：h=here, t=top, b=bottom, p=獨立頁
%     \centering                   — 圖片置中
%     \includegraphics[選項]{檔名} — 載入圖片
%       檔名：可帶副檔名 (head.jpg) 或不帶 (head)
%             不帶時 LaTeX 自動搜尋 .pdf → .png → .jpg
%             搜尋路徑已設定為 figures/ 和 assets/（見 preamble.tex）
%       常用選項：
%         width=0.8\linewidth      — 寬度（\linewidth=當前欄寬, \textwidth=頁寬）
%         height=5cm               — 高度
%         keepaspectratio           — 同時限制寬高時保持比例
%     \caption{說明文字}           — 圖片標題（自動編號）
%     \label{fig:xxx}              — 引用標籤（用 \cref{fig:xxx} 自動產生 "Figure 1"）
%   \end{figure}

\subsection{Basic Usage}

\begin{figure}[htbp]
    \centering
    \includegraphics[width=\linewidth]{head}
    \caption{範例圖片（預設滿版寬）。}
    \label{fig:demo-full}
\end{figure}

Reference: \cref{fig:demo-full} shows the full-width figure.

\subsection{Size Control}

% 限制寬度
\begin{figure}[htbp]
    \centering
    \includegraphics[width=0.3\linewidth]{head}
    \caption{縮到 30\% 頁寬。}
    \label{fig:demo-30}
\end{figure}

% 限制高度
\begin{figure}[htbp]
    \centering
    \includegraphics[height=4cm]{head}
    \caption{限制高度 4cm。}
    \label{fig:demo-4cm}
\end{figure}

% 同時限制寬高，保持比例（不超過任一限制）
% \includegraphics[width=0.5\linewidth, height=3cm, keepaspectratio]{head}

\subsection{Side-by-Side Figures}

% 兩張圖水平並排，各有自己的 caption 和 label
% minipage 寬度加起來要 < 1\linewidth（留一點給 \hfill 的間距）
\begin{figure}[htbp]
    \centering
    \begin{minipage}[t]{0.48\linewidth}
        \centering
        \includegraphics[width=\linewidth]{head}
        \caption{左圖。}
        \label{fig:side-left}
    \end{minipage}
    \hfill
    \begin{minipage}[t]{0.48\linewidth}
        \centering
        \includegraphics[width=\linewidth]{head}
        \caption{右圖。}
        \label{fig:side-right}
    \end{minipage}
\end{figure}

See \cref{fig:side-left} and \cref{fig:side-right}.

\subsection{Shadow Box}

% \shadowbox 是定義在 macros.tex 的通用包裝器，加上邊框 + 模糊陰影
% 不限於圖片，任何內容都能包：\shadowbox{文字或 box}
% 底層使用 tcolorbox 的 fuzzy shadow，需要 \usetikzlibrary{shadows}（見 preamble.tex）

\begin{figure}[htbp]
    \centering
    \shadowbox{\includegraphics[width=0.2\linewidth]{head}}
    \caption{圖片加上邊框與模糊陰影（\texttt{\textbackslash shadowbox}）。}
    \label{fig:demo-shadow}
\end{figure}

% ============================================================
\section{Tables}
% ============================================================

% table 語法解析：
%   \begin{table}[htbp]       — 浮動環境，位置偏好同 figure
%   \centering                 — 置中
%   \caption{} + \label{}      — 標題和引用標籤
%   \begin{tabular}{欄位定義}  — 表格內容
%     欄位對齊：l=靠左, c=置中, r=靠右（幾個字母 = 幾欄）
%     &                        — 分隔欄位
%     \\                       — 換行
%     \toprule                 — 頂線（booktabs 提供，比 \hline 好看）
%     \midrule                 — 中線（表頭與內容之間）
%     \bottomrule              — 底線
%   \end{tabular}

\begin{table}[htbp]
    \centering
    \caption{A sample table with booktabs.}
    \label{tab:sample}
    \begin{tabular}{lcc}
        \toprule
        \textbf{Name} & \textbf{Score} & \textbf{Grade} \\
        \midrule
        Alice   & 95 & A+ \\
        Bob     & 87 & A  \\
        Charlie & 72 & B  \\
        \bottomrule
    \end{tabular}
\end{table}

See \cref{tab:sample}.

% ============================================================
\section{Algorithm}
% ============================================================

% algorithm2e 語法解析：
%   \begin{algorithm}[H]          — [H] 強制放在這裡（不浮動）
%   \caption{} + \label{}         — 標題和引用標籤
%   \KwInput{}, \KwOutput{}       — 輸入輸出描述
%   \;                            — 行尾分號（每行結尾都要加）
%   \While{條件}{程式碼}          — while 迴圈
%   \For{條件}{程式碼}            — for 迴圈
%   \uIf{條件}{} \uElseIf{}{} \Else{} — if-elseif-else（u 表示不加結尾橫線）
%   \Return{值}                   — 回傳
%   數學模式：直接用 $ ... $ 寫變數和運算式

\begin{algorithm}[H]
    \caption{Binary Search}
    \label{alg:bsearch}
    \KwInput{Sorted array $A[1..n]$, target $t$}
    \KwOutput{Index of $t$ in $A$, or $-1$}
    $lo \gets 1$, $hi \gets n$\;
    \While{$lo \le hi$}{
        $mid \gets \lfloor (lo + hi) / 2 \rfloor$\;
        \uIf{$A[mid] = t$}{
            \Return{$mid$}\;
        }
        \uElseIf{$A[mid] < t$}{
            $lo \gets mid + 1$\;
        }
        \Else{
            $hi \gets mid - 1$\;
        }
    }
    \Return{$-1$}\;
\end{algorithm}

See \cref{alg:bsearch}.

% ============================================================
\section{Cross-References}
% ============================================================

% \cref 由 cleveref 套件提供，自動產生類型名稱 + 編號
% 根據 label 前綴自動判斷類型：
%   \label{fig:xxx}  → \cref{fig:xxx}  → "Figure 1"
%   \label{tab:xxx}  → \cref{tab:xxx}  → "Table 1"
%   \label{eq:xxx}   → \cref{eq:xxx}   → "Equation (1)"
%   \label{alg:xxx}  → \cref{alg:xxx}  → "Algorithm 1"
%   \label{sec:xxx}  → \cref{sec:xxx}  → "Section 1"
%
% 一次引用多個：\cref{fig:side-left,fig:side-right} → "Figures 1 and 2"
%
% \Cref 用於句首（大寫開頭）：\Cref{fig:xxx} → "Figure 1"
% \cref 用於句中（小寫開頭）：\cref{fig:xxx} → "figure 1"（視語言設定）

See \cref{fig:demo-full}, \cref{tab:sample}, \cref{eq:gaussian}, and \cref{alg:bsearch}.

% ============================================================
\section{Footnotes and Margin Notes}
% ============================================================

% --- 腳註 ---
% \footnote{內容} — 在當前位置加上編號，頁尾顯示內容
This sentence has a footnote\footnote{This is the footnote content, displayed at the bottom of the page.}.

% --- 邊欄筆記 ---
% \marginmark{內容} — 在頁面邊欄顯示灰色小字（自訂指令，定義在 macros.tex）
This paragraph has a margin note.\marginmark{Remember this point.}

% ============================================================
\section{Colors}
% ============================================================

% --- 文字變色 ---
% \textcolor{顏色}{文字}
\textcolor{red}{Red}, \textcolor{blue}{Blue}, \textcolor{ForestGreen}{ForestGreen}.

% --- 整段變色 ---
% {\color{顏色} 這段文字都是這個顏色}
{\color{BrickRed} This entire paragraph is in BrickRed. The color applies until the closing brace.}

% --- 自訂顏色 ---
% \definecolor{名稱}{模式}{值}
%   模式 HTML: \definecolor{MyOrange}{HTML}{FF6B35}
%   模式 RGB:  \definecolor{MyBlue}{rgb}{0.1, 0.3, 0.8}（0~1 範圍）
%
% --- 顏色的 ! 混合語法 ---
% A!比例!B = 取比例% 的 A + 剩下的 B
%   blue!30!white  = 30% 藍 + 70% 白 → 淡藍
%   red!50!black   = 50% 紅 + 50% 黑 → 暗紅
%   gray!10        = gray!10!white 的簡寫 → 10% 灰 + 90% 白 = 非常淺的灰
% 可以連鎖：red!50!blue!30!white = 先混出 50%紅+50%藍，再取 30% 混白
\definecolor{CustomPurple}{HTML}{7B2D8E}
\textcolor{CustomPurple}{Custom purple text}. Also: \textcolor{blue!50!black}{50\% blue + 50\% black}.

% --- 背景高亮 ---
% \colorbox{背景色}{文字}
\colorbox{yellow!30}{highlighted text} in a sentence.

% ============================================================
\section{Multi-Column Layout}
% ============================================================

% twocol / threecol — 兩欄 / 三欄排版
% \columnbreak       — 強制在此處換到下一欄

\subsection{Two Columns}

\begin{twocol}
\textbf{Left column.} Lorem ipsum dolor sit amet, consectetur adipiscing elit. Sed do eiusmod tempor incididunt ut labore.

\columnbreak

\textbf{Right column.} Ut enim ad minim veniam, quis nostrud exercitation ullamco laboris nisi ut aliquip.
\end{twocol}

\subsection{Three Columns}

\begin{threecol}
\textbf{Column 1.} Short text.

\columnbreak

\textbf{Column 2.} Medium length paragraph for demonstration.

\columnbreak

\textbf{Column 3.} Another column.
\end{threecol}

% ============================================================
\section{Page Control}
% ============================================================

% \newpage     — 強制換頁
% \clearpage   — 強制換頁，並輸出所有尚未放置的浮動環境（figure, table）
%                如果有圖表被 LaTeX 延遲放置，\clearpage 會先把它們全部輸出
% \pagebreak   — 建議換頁（LaTeX 可能忽略）
% \nopagebreak — 建議不要在此換頁

% 一般用 \clearpage 比 \newpage 安全，因為它確保浮動環境不會跑到錯誤的位置

% ============================================================
\section{Todo Examples}
% ============================================================

% \todox{內容}     — 黃色行內提醒（顯眼的備忘）
% \todomark{內容}  — ☐ 未完成項目
% \tododone{內容}  — ☒ 已完成項目（文字加刪除線）

\todox{This is an inline todo note for review.}

Task list:
\begin{itemize}[label={}]
    \item \todomark{Write the introduction}
    \item \todomark{Add figures}
    \item \tododone{Set up template}
\end{itemize}

% ============================================================
\section{Bibliography}
% ============================================================

% 1. 建立 .bib 檔案（例如 references.bib），內容格式：
%    @article{key,
%        author = {Author Name},
%        title  = {Paper Title},
%        journal = {Journal Name},
%        year   = {2024},
%    }
%    可從 Google Scholar 直接匯出 BibTeX 格式
%
% 2. 在 preamble.tex 的 biblatex 設定下方加入：
%    \addbibresource{references.bib}
%
% 3. 在文中引用：\cite{key} → [1]
%    或 \textcite{key} → "Author [1]"
%
% 4. 在文末放置參考文獻列表：
%    \printbibliography
%
% 5. 編譯順序：lualatex → biber → lualatex → lualatex

% \printbibliography

\end{document}
